\begin{table}[H]
\color{blue}\centering\footnotesize
\setlength{\tabcolsep}{4pt}
\caption{네 명의 공존을 평가하기 위한 지표 후보. $K_j$는 판 $j$의 완주자 수이고 $J$는 판 수다. $c_{ij}$는 자리 $i$가 완주하면 1, 아니면 0이다.}\label{tab:coexistence-candidates}
\begin{tabularx}{\textwidth}{@{}p{.16\textwidth}p{.26\textwidth}X@{}}
\toprule
지표와 용도 & 정의 & 답하는 질문과 해석 범위 \\
\midrule
네 명 완주율(주 지표) & $J^{-1}\sum_j\mathbf{1}\{K_j=4\}$ & 네 명 모두 끝까지 남은 판이 얼마나 많은가. 세 명 완주는 성공에 포함하지 않는다. \\
개인 완주 비율(보조) & $J^{-1}\sum_j K_j/4$ & 전체 자리 중 얼마나 많이 완주했는가. 한 판에 세 명만 남아도 75\%이므로 네 명 공존과 구별한다. \\
완주자 수별 누적 비율(보조) & $J^{-1}\sum_j\mathbf{1}\{K_j\geq q\}$, $q=1,2,3,4$ & 최소 몇 명까지 함께 남았는가. 네 기준을 함께 보고하며 $q=4$가 주 지표와 같다. \\
쌍 완주 비율(보조) & $J^{-1}\sum_j K_j(K_j-1)/12$ & 같은 판의 두 자리가 함께 완주한 비율은 얼마인가. 세 명 완주도 50\%이므로 네 명 완주를 대신하지 못한다. \\
완주 여부의 기하평균 & $J^{-1}\sum_j(\prod_{i=1}^4 c_{ij})^{1/4}$ & 한 자리라도 완주하지 못하면 0이다. 이진 완주 여부에서는 네 명 완주율과 같아 별도 지표로 추가하지 않는다. \\
Sen 지수(자원 분배 보조) & 판별 최종 토큰 평균 $\times(1-G_j)$의 평균 & 토큰이 얼마나 많고 고르게 남았는가. 상금·EXIT 보너스·잔량을 반영하므로 완주율을 대신하지 않는다. \\
\bottomrule
\end{tabularx}
\end{table}
